\section{The problem, results, and proof strategy}\label{ms:introduction}
\subsection{The problem and the candidate constant}
Let $\mathcal P$ be the laws on $\R$ with mean zero, variance one and
finite third absolute moment $\rho(P)=\E|X|^3$. For independent copies
$X_1,\ldots,X_n$ of $P$, put $S_n=\sum_iX_i$ and
\[
 \Delta_n(P)=\sup_{H\in\R}|\Pp_P(S_n\le H)-\Phi(H/\sqrt n)|,
 \qquad C_n=\sup_{P\in\mathcal P}\frac{\sqrt n\,\Delta_n(P)}{\rho(P)}.
\]
Here $\Phi$ and $\phi$ are the standard normal distribution function
and density. The sharp-constant conjecture is
\begin{equation}\label{ms:conjecture}
 \sup_{n\ge1}C_n=\CE,
 \qquad \CE=\frac{3+\sqrt{10}}{6\sqrt{2\pi}}=0.409732\ldots.
\end{equation}
Throughout, $H$ is a threshold for the unnormalized sum $S_n$, and
a terminating decimal in an inequality denotes the corresponding exact
rational number. Affine standardization gives the equivalent result
for any mean and positive variance, with third absolute central moment
divided by the cube of the standard deviation.

Berry and Esseen established the classical order $n^{-1/2}$
\cite{Berry1941,Esseen1942,Esseen1945}. Esseen's lower bound
\cite{Esseen1956} explains the proposed constant. If
$p_*=(4-\sqrt{10})/2$, $q_*=1-p_*$ and
$X=(B-p_*)/\sqrt{p_*q_*}$ for $B\sim\operatorname{Bernoulli}(p_*)$,
the fixed-law lattice limit gives
\begin{equation}\label{ms:bernoulli-limit}
 \lim_{n\to\infty}\frac{\sqrt n\Delta_n(P)}{\rho(P)}
 =\frac{3+q_*-p_*}{6\sqrt{2\pi}(p_*^2+q_*^2)}=\CE.
\end{equation}
Indeed $p_*^2+q_*^2=10-3\sqrt{10}$ and
$(3+q_*-p_*)/(p_*^2+q_*^2)=3+\sqrt{10}$.
Every constant below $\CE$ therefore fails for a sufficiently large sum
of this one law. The upper conjecture is associated with Zolotarev
\cite{Zolotarev1966,Zolotarev1967}; Schulz proves the strict finite-index
bound for every nondegenerate Bernoulli law \cite{Schulz2016}.


\subsection{Results and their hypotheses}\label{ms:scope}
The results have three distinct roles. We state their interfaces here;
the cited theorems give their complete hypotheses and proofs.
\begin{center}\small
\begin{tabular}{@{}P{.25\linewidth}P{.68\linewidth}@{}}
\toprule
Result & Conclusion and required inputs \\
\midrule
Local comparison, Theorem~\ref{short:count:local} &
For $n\ge700$, every law in its specified conditional-support region,
with $p\in[3/8,93/200]$, $0<\lambda\le7/20$ and $|H|\le46/25$,
has excess below $-\lambda/100$. The zero-noise excess is strictly
negative. This statement uses only its explicit law hypotheses. \\
Localization, Theorem~\ref{v9:entry:theorem} &
A maximizer with nonnegative excess at $n\ge42$ satisfies
$nD\le32.20665378$ and $|H|\le1.83605680$, under
$C_{n-1}\le\CE$, the structural inequality and the two stated
finite-verification conditions. \\
Global implication, Theorem~\ref{short:global} &
With the retained structural, binomial and finite-verification inputs,
a base $C_j\le\CE$ for every $j\le699$ implies $C_n<\CE$ for
$n\ge700$ and $\sup_n C_n=\CE$. The finite base remains open here. \\
\bottomrule
\end{tabular}
\end{center}
The structural input is displayed in \eqref{v9:entry:structural}; the
finite-binomial input is \eqref{short:count:binomial}.
Conditions~\ref{v9:entry:two-pass-condition} and
\ref{v9:entry:signed-condition} specify the finite localization
checks. Their exact replay is recorded in Appendix~\ref{revision:reproducibility}.
The analytical reduction, the whole-domain coverage, and the
correspondence of the checking code with those inequalities are the
components of this computational proof. We retain these premises
explicitly in the theorem statements.

\subsection{Why the defect controls a maximizer}
At a fixed $n$, maximize the difference between the two sides of the
candidate inequality:
\begin{equation}\label{v9:intro:excess}
 a(P,H)=\sqrt n\{\Pp_P(S_n\le H)-\Phi(H/\sqrt n)\}-\CE\rho(P).
\end{equation}
Proposition~\ref{v9:core:attainment} proves attainment. Its first-order
condition bounds $y\mapsto\Pp(S_{n-1}\le H-y)$ by an explicit
piecewise cubic function that agrees with it on the closed support.
This contact condition allows us to use the geometry of that support.

Let $\tau(t)=\E[X\ind_{\{X\ge t\}}]$ be the zero-bias density and
$J=\sum_g\tau_g|g|^2$, where $g$ runs over bounded support gaps and
$\tau_g$ is the constant value of $\tau$ there. A lattice law of span
$h$ has $J\ge h$; if every successive gap has length $h$, then $J=h$.
Thus $J$ measures gap size with zero-bias weights. The inequality
\[
 \gamma+3J\le(3+\sqrt{10})\rho,\qquad \gamma=\E X^3,
\]
proved in Lemma~\ref{ms:physical-defect}, extends the moment--span
bound from lattice laws to arbitrary standardized laws. Its slack
\[
 D=(3+\sqrt{10})\rho-\gamma-3J
\]
is nonnegative and vanishes exactly at the positive-skew Bernoulli
law in \eqref{ms:bernoulli-limit}. The lattice consequence and its
reflection are proved in Corollary~\ref{opus:lattice-gap}.

Integrating the contact relation against $\tau$ in two ways gives the
conservation identity of Theorem~\ref{revision:contact-conservation}.
Since $\CE=(3+\sqrt{10})\phi(0)/6$, its two leading moment terms
combine into a multiple of $D$. Solving for the excess gives
\[
 a=\frac{\mathcal R_c-n\CE\rho(1-e)
             -n\phi(0)eD/6-n^{3/2}L}{n+H^2},
 \qquad e=\exp(-H^2/(2n)).
\]
Here $L\ge0$ is a weighted contribution from the law of $H-S_{n-1}$
inside the support gaps, and $\mathcal R_c$ is an explicitly defined
remainder. The three subtracted terms are nonnegative. A nonnegative
excess therefore requires $n\phi(0)eD/6\le\mathcal R_c$.
Once localization bounds both $\mathcal R_c$ and $|H|$, this bounds
$nD$. The uniform bounds come from Theorem~\ref{v9:entry:theorem},
not from a fixed-law asymptotic expansion.

The predecessor hypothesis supplies an upper bound for
$\Pp(S_{n-1}\le H-t)$ in the contact condition, which excludes remote
support points; see \eqref{v9:entry:contact}. This support argument is
used both in coarse localization and in the final two-ray estimate.
With $nD\le32.20665378$, the geometry at $n\ge700$ gives
$3/8<p<93/200$, inward support displacement below $1/250$ and
$\lambda<7/20$. Count conditioning then compares the law with a
binomial sum. One threshold count controls noise up to $1/100$;
five neighboring counts control the remaining interval. Both
comparisons keep all probabilities entering from other counts.

\begin{figure}[ht]
\centering
\begin{tikzpicture}[node distance=4mm,>=Latex,
 box/.style={draw,rounded corners,align=center,text width=.86\linewidth,
 inner sep=5pt,font=\small}]
\node[box] (contact) {Attainment and contact; nonnegative gap defect $D$};
\node[box,below=of contact] (identity) {Zero-bias conservation: an upper bound for the remainder controls $nD$};
\node[box,fill=gray!10,below=of identity] (entry) {Structural inequality + predecessor support bound + finite verification\\
Coarse localization: $nD\le32.20665378$, $|H|\le1.83605680$};
\node[box,below=of entry] (local) {For $n\ge700$: geometry and support bounds place the maximizer\\
in the region of the local binomial comparison; the excess is negative};
\node[box,line width=.8pt,below=of local] (global) {Finite base through $699$ (unproved here) + least-failure argument\\
Conditional global sharp inequality};
\draw[->] (contact)--(identity);
\draw[->] (identity)--(entry);
\draw[->] (entry)--(local);
\draw[->] (local)--(global);
\end{tikzpicture}
\caption{The proof's dependencies. The shaded step uses the specified
finite verification; the final step also requires the unproved finite base.
The local comparison is a theorem for its explicit law region, independently
of the preceding localization.}
\label{opus:proof-flow}
\end{figure}

\FloatBarrier
\subsection{The computer-assisted localization}
Analysis restricts a possible maximizer to a bounded moment region.
It also bounds the remainder by functions of moment and support
parameters. The computation encloses these functions on finitely many
closed boxes. Each box is either excluded by a necessary inequality,
or assigned bounds valid throughout that box. The boxes cover their
parent regions, including shared boundaries, so every possible
maximizer retains a valid bound.

The certificate is a compact record of this process. It stores
checkpoints for $24$ third-moment cells over four signed-moment passes,
including support radii, bounds and descendant counts; the individual
$11264$ nominal boxes are reconstructed. Two preliminary coarse passes
are reconstructed as well. The exporter chooses a successful sequence
of passes and rounds its reported bounds outward. The acceptor
recomputes the prescribed subdivisions and verifies their inequalities
and output bounds.

All numerical acceptance uses integers and exact rational numbers,
with directed rational enclosures for roots, exponentials and Gaussian
constants. The exploratory program and the checker share the routine
\texttt{signed\_core.stage}; acceptance is an exact replay of that
calculation. Independent checking requires auditing these routines
against the mathematical rules or implementing those rules separately.
Appendix~\ref{revision:reproducibility} gives a code-reading guide,
reproduction commands and the recorded execution environment.

Subdivision and rigorous enclosure have established precedents in
computer-assisted proofs \cite{SolovyevHales2013}. Numerical optimization
also appears in Berry--Esseen estimates \cite[Section~2]{Shevtsova2011}.
The custom JSON format organizes this calculation; the rigor comes from
the valid enclosures, coverage and correct implementation of the rules.

\subsection{Related work}
Fourier and structural moment inequalities progressively improved the
universal upper bound \cite{vanBeek1972,Prawitz1972,Tyurin2010,
Shevtsova2011,Shevtsova2013}. We use Shevtsova's structural inequality
explicitly in localization. Recent manuscripts report smaller upper
constants \cite{xiao2026berryesseen,lin2026berryesseen}, and He and Cheng
report eventual sharpness beyond a uniform, very large threshold
\cite{hecheng2026eventually}. Work on independent arrays and extremizers
addresses the related non-iid problem \cite{BellIndependent2026,
BellExtremizers2026}. These recent reported results are not used in
our proof. An earlier AI-assisted working draft is on Hexagon
\cite{rodziewiczkolodziejek2026momentdefects}.
